\section{Praktikum}
\label{sec:praktikum-bab12}

\begin{praktikumbox}
\vspace{1.1em}
\textbf{Judul:} Membandingkan mode pengalamatan pada akses array.\\[0.3em]
\textbf{Sub-CPMK:} 3.3 \quad\textbullet\quad
\textbf{Alokasi waktu:} 100 menit \quad\textbullet\quad
\textbf{Bentuk:} individu
\end{praktikumbox}

\subsection{Tujuan Praktikum}

Setelah menyelesaikan praktikum ini, mahasiswa mampu menggunakan
tujuh mode pengalamatan, menghitung alamat efektif, dan membuktikan
bahwa beberapa mode dapat menunjuk elemen array yang sama.

\subsection{Alat dan Bahan}

\begin{itemize}
    \item Komputer dengan emu8086 terpasang.
    \item DOSBox dengan TASM, TLINK, dan DEBUG.
    \item Berkas \file{mode-array.asm} dari \file{code/asm/}.
    \item Lembar pengamatan register dan alamat.
\end{itemize}

\subsection{Langkah Kerja}

\begin{enumerate}
    \item \textbf{Merakit dan menjalankan program.}
    \begin{enumerate}
        \item Buka DOSBox, lalu masuk ke direktori kerja.
        \item Rakit dengan \file{tasm mode-array.asm}.
        \item Link dengan \file{tlink mode-array.obj}.
        \item Jalankan program dan catat keluarannya.
    \end{enumerate}

    \item \textbf{Menelusuri tiga mode.}
    \begin{enumerate}
        \item Buka program pada emu8086 dan jalankan
              \eng{single step}.
        \item Catat \reg{AX}, \reg{SI}, dan \reg{BX} pada bagian
              \textit{direct}, \textit{direct index}, dan
              \textit{base relative}.
        \item Hitung alamat efektif setiap akses dengan rumus
              $EA = basis + indeks + displacement$.
        \item Bandingkan nilai \code{jumlah1}, \code{jumlah2}, dan
              \code{jumlah3} pada jendela memori.
    \end{enumerate}

    \item \textbf{Mengubah isi array.}
    \begin{enumerate}
        \item Ubah elemen array menjadi \code{10, 20, 30, 40, 100}.
        \item Rakit ulang dan jalankan program.
        \item Pastikan ketiga cara tetap menghasilkan jumlah yang
              sama.
        \item Jelaskan mengapa perubahan isi tidak mengubah
              hubungan antarmode.
    \end{enumerate}

    \item \textbf{Mencoba mode lain.}
    \begin{enumerate}
        \item Ubah cara kedua dari \lstasm|data1[SI]| menjadi
              \lstasm|[SI+data1]|.
        \item Ubah cara ketiga dari \lstasm|[BX+SI]| menjadi
              \lstasm|[BX+4]| untuk elemen ke-2.
        \item Rakit ulang dan pastikan hasilnya tetap sama.
    \end{enumerate}

    \item \textbf{Memeriksa alamat dengan DEBUG.}
    \begin{enumerate}
        \item Jalankan \file{debug mode-array.exe}.
        \item Catat nilai \reg{DS} dengan perintah \file{r}.
        \item Hitung alamat fisik awal array dengan rumus
              $(\reg{DS} \times \hex{10}) + \text{offset data1}$.
        \item Tampilkan isi array dengan
              \lstasm|D DS:offset L 0A|.
        \item Cocokkan urutan word dengan nilai array.
    \end{enumerate}

    \item \textbf{Menulis potongan tujuh mode.}
    \begin{enumerate}
        \item Tulis potongan program yang mengakses elemen ke-2
              dengan mode \textit{direct}, \textit{indirect},
              \textit{base relative}, \textit{direct index}, dan
              \textit{based indexed}.
        \item Jelaskan mode mana yang tidak mengakses memori.
    \end{enumerate}
\end{enumerate}

\subsection{Hasil yang Dikumpulkan}

\begin{enumerate}
    \item Berkas \file{mode-array.asm} versi modifikasi.
    \item Tangkapan layar keluaran program.
    \item Tabel alamat efektif untuk tiga mode.
    \item Tangkapan layar isi array pada DEBUG.
    \item Penjelasan perbedaan mode \textit{indirect} dan
          \textit{direct index}.
\end{enumerate}

\begin{asesmenbox}
\vspace{1.1em}
\textbf{Asesmen praktikum 10 (bobot 5\% dari komponen praktikum).}

\begin{tabular}{@{}p{6.2cm}cp{4.6cm}@{}}
    \toprule
    \textbf{Indikator} & \textbf{Bobot} & \textbf{Kriteria} \\
    \midrule
    Ketepatan penggunaan tiga mode & 35\% &
    Ketiga cara mengakses elemen dengan benar \\
    \addlinespace
    Ketepatan perhitungan alamat efektif & 30\% &
    Basis, indeks, dan displacement dijumlahkan benar \\
    \addlinespace
    Ketepatan analisis hasil & 20\% &
    Kesamaan hasil dijelaskan melalui alamat efektif \\
    \addlinespace
    Kerapian dokumentasi & 15\% &
    Lembar kerja lengkap dan dapat diulang \\
    \bottomrule
\end{tabular}
\end{asesmenbox}
